\notechapter{N-20230314}{Category Theory Continued: Free and Forgetful Functors, Concrete Categories, and Group Actions}

\section{Free and forgetful functors}

The first page discusses free and forgetful functors.  These words are often used informally before one learns the general adjunction language.  The basic pattern is simple: a forgetful functor removes structure, while a free functor freely adds structure.

For example, the forgetful functor
\[
  U:\mathbf{Vect}_k\to\mathbf{Set}
\]
sends a vector space to its underlying set and a linear map to its underlying function.  Similarly,
\[
  \mathbf{Grp}\to\mathbf{Set},\qquad
  \mathbf{CRing}\to\mathbf{Set}
\]
forget group or ring structure.  There is also a forgetful functor
\[
  \mathbf{CRing}\to\mathbf{Grp}
\]
that sends a commutative ring $R$ to its additive abelian group $(R,+)$.

A free functor goes in the opposite direction.  The free vector space functor
\[
  F:\mathbf{Set}\to\mathbf{Vect}_k
\]
sends a set $S$ to the vector space with basis $S$.  A function $f:S\to T$ induces a linear map $F(S)\to F(T)$ by sending the basis vector $s$ to the basis vector $f(s)$ and extending linearly.

The slogan is:
\[
\begin{gathered}
  \text{free} = \text{add the least structure needed},\\
  \text{forgetful} = \text{discard structure but keep underlying data}.
\end{gathered}
\]

\section{Faithful and full functors}

For a functor $F:\mathcal A\to\mathcal B$, each pair of objects $A_1,A_2$ gives a map on hom-sets:
\[
  F:\operatorname{Hom}_{\mathcal A}(A_1,A_2)
  \to
  \operatorname{Hom}_{\mathcal B}(F(A_1),F(A_2)).
\]
The functor is faithful if this map is injective for every pair of objects.  It is full if this map is surjective for every pair.

Faithful means that distinct arrows remain distinct after applying $F$.  Full means that every arrow between the images comes from an arrow upstairs.  A forgetful functor is often faithful but not full.  For instance, not every function between underlying sets of groups is a group homomorphism.

\section{Concrete categories}

A concrete category is a category equipped with a faithful functor to $\mathbf{Set}$:
\[
  U:\mathcal A\to\mathbf{Set}.
\]
This means every object has an underlying set and every arrow has an underlying function, in a way that does not collapse distinct arrows.

The handwritten note is careful here: many familiar categories are concrete, such as groups, rings, vector spaces, and topological spaces.  But the word “concrete” is extra structure, not just a vibe.  It depends on choosing a faithful functor to sets.

\section{Functors out of a one-object group category}

Let $G$ be a group and let $\mathcal G$ be its associated one-object category.  There is one object $*$, and
\[
  \operatorname{Hom}_{\mathcal G}(*,*)=G.
\]
Consider a functor
\[
  F:\mathcal G\to\mathbf{Set}.
\]
The object $*$ is sent to a set
\[
  S=F(*).
\]
Each group element $g\in G$, regarded as an arrow $*:\to*$, is sent to a function
\[
  F(g):S\to S.
\]
Since functors preserve identity and composition,
\[
  F(e)=\operatorname{id}_S,\qquad
  F(gh)=F(g)\circ F(h)
\]
up to the convention of left or right actions.  This is exactly a group action of $G$ on $S$.

Thus:
\[
  \text{functor }\mathcal G\to\mathbf{Set}
  \quad\Longleftrightarrow\quad
  \text{a }G\text{-set}.
\]
This is an early example of category theory turning algebra into actions that preserve structure.

\section{Functors between one-object group categories}

If $G$ and $H$ are groups and $\mathcal G,\mathcal H$ their one-object categories, then a functor
\[
  F:\mathcal G\to\mathcal H
\]
must send the unique object of $\mathcal G$ to the unique object of $\mathcal H$.  On arrows, it gives a function
\[
  G\to H.
\]
Functoriality requires
\[
  F(gh)=F(g)F(h),\qquad F(e_G)=e_H.
\]
So functors $\mathcal G\to\mathcal H$ are exactly group homomorphisms $G\to H$.

This shows again that functoriality is not an extra decoration.  It is exactly the condition that algebraic composition is preserved.

\section{Why the note says “learn algebra first”}

The handwritten page remarks that one should learn enough algebra before pushing further.  That is reasonable: the example of a group as a one-object category is conceptually simple only if one already understands group operations, homomorphisms, and group actions.  Without those, the category-theoretic formulation can feel like empty abstraction.

The actual learning path is circular in a productive way.  Algebra supplies examples for category theory; category theory then reorganizes algebra.

\section{Looking ahead: natural transformations}

The final handwritten comment says that natural transformations are coming next.  This is the right next step.  If functors are structure-preserving maps between categories, then natural transformations are structure-preserving maps between functors.

In the group-action example, a natural transformation between two functors $\mathcal G\to\mathbf{Set}$ is a function between the underlying $G$-sets that commutes with the $G$-action.  In other words, it is an equivariant map.  This is the categorical reason natural transformations are not optional: they are the morphisms between representations of the same structure.
